
Rebuttal (Reviewer 1)

We thank the reviewer again. The first message covered Q1--Q6. Below is the continuation (Q7--Q21).

Q7 (Write the loss and make the gradient-flow structure explicit):

We thank the reviewer for this suggestion and will make the statement explicit. We will define the population loss
\[
\mathcal L(W)=\mathbb E_Z[\mathcal L(Y,\hat y(X;W))],\qquad
d\mathcal L(Z;W)=\partial_{\hat y}\mathcal L(Y,\hat y(X;W)),
\]
and state that the mean-field ODE is the gradient flow (with learning-rate schedules $\xi_\ell(t)$) of $\mathcal L$ in the mean-field parameterization. These definitions are currently spelled out in the appendix; for the camera-ready, we will also add them in the main body (with concise notation) so Section 3 is self-contained.

Q8 (Learning-rate schedules):

We will define $\xi_\ell(t)\ge 0$ as the learning-rate prefactor in the mean-field ODE (allowing time-dependent schedules such as constant, piecewise, exponential, or cosine forms). We will also state the discrete-time correspondence explicitly: with step size $\varepsilon$ and iteration index $n$, Euler discretization gives
\[
W_{n+1}=W_n-\varepsilon\,\xi_\ell(\varepsilon n)\,(\text{drift at time }\varepsilon n),
\]
and the mean-field limit is $\varepsilon\to 0$ with $t=\varepsilon n$ (plus the usual width scaling). In the classical mean-field parameterization, this corresponds to the familiar $1/n$ learning-rate scaling at width $n$.

Q9 ("Better than null" assumption):

We will restate this as a non-degeneracy condition: the initial predictor should perform strictly better than the trivial constant predictor (e.g. $\hat y\equiv \phi_L(0)$). In the manuscript appendix, this appears as a sufficient condition ensuring the limit point is not identically zero (the Non-Degeneracy assumption: $\mathcal L(W(0))<\mathbb E[\mathcal L(Y,\phi_L(0))]$).

Q10 (ReLU vs. $\phi_2'$ bounded away from zero; high probability in $r$):

We will clarify both points.

- Meaning of $r$. $r$ is the rank (number of channels) in the low-rank bottleneck. The second-layer pre-activation is a sum over $r$ channels:
\[
H_2(c_2;x,W)=\sum_{k=1}^r W_{c_2,k}f_k(x;W).
\]

- Why ReLU is excluded in the formal global-convergence proof. Our regularity assumption in the manuscript appendix requires $\inf_u |\phi_2'(u)|>0$, which excludes plain ReLU (since $\phi'(u)=0$ for $u\le 0$). This is used in the step "vanishing conditional gradient implies vanishing residual".

- What we can still say for ReLU (high-level). The only problematic case is convergence to a degenerate stationary point where all gates are off ($H_2\le 0$ almost everywhere). In our architecture, this corresponds to a rare alignment event across $r$ channels; we treat the resulting high-probability-in-$r$ claim as intuition, not as a theorem.

To make the logical use of the nonzero-derivative condition explicit, the global-optimality proof uses the implication
\[
\mathbb E[d\mathcal L(Z;W)\cdot\phi_2'(H_2(c_2;X,W))\mid X=x]=0
\Rightarrow
\mathbb E[d\mathcal L(Z;W)\mid X=x]=0,
\]
which holds when $\phi_2'$ is bounded away from $0$ on the relevant set. For ReLU, the failure mode is the dead-gate event: $\phi_2'(H_2)=0$ even though $d\mathcal L\neq 0$, so the conditional moment can vanish at a non-optimal point.

In RF-LR, $H_2(c_2;x,W)=\sum_{k=1}^r W_{c_2,k}f_k(x;W)$ is a sum over $r$ channels. A design intuition is that sign-structured/NMF-style initialization can bias $H_2$ positive on target regions and reduce dead-gate behavior (not a proof step). Under symmetric random mixing (e.g., Xavier/Gaussian draws for $W_{c_2,k}$) and non-degenerate $f_k(x)$, one-neuron off probability is approximately $1/2$; therefore the event that all channels are simultaneously dead is heuristically exponentially unlikely in $r$ (informally $e^{-O(r)}$), i.e., concentrated on an input subset of exponentially small measure. We keep this as intuition, not theorem (GeLU satisfies assumptions directly).

We will also add a precise definition of mixing and the seminorm $\|W\|_{\infty,1}$.

Q11 (What is mixing?):

We will define mixing explicitly: $W^{(\ell)}$ recombines $r$ channels into neuron pre-activations, with structural bound $\|W^{(\ell)}\|_{\infty,1}\le rK$. We will also state that $W^{(\ell)}$ may be drawn from any bounded-support distribution (or any law satisfying the same almost-sure bound).

Q12 (Lipschitz notation in Sec. 3.5--3.6):

We agree the notation was confusing. In this part, the trainable object is factor $A$, while the prefactor is a norm of the frozen mixing matrix. We will correct the inequality, explicitly separate trainable vs.\ frozen quantities, and harmonize notation in the camera-ready to avoid switching between $W$/$L$ symbols for the same mixing entries.

Q13 (Solution operator $F$):

We will define $F$ as the Picard integral map whose fixed point is the ODE solution trajectory, and state this clearly where first introduced.

Q14 (Minor punctuation):

We will correct the sentence punctuation and improve readability.

Q15 (Theorem~3.2 scope in depth/width):

Thank you for this clarification request. The result is intended to hold for arbitrary depth in our RF-LR architecture, with the multilayer extension detailed in the appendix. In the camera-ready version, we will rewrite the theorem statement to specify explicitly which widths tend to infinity (all hidden-layer widths under mean-field scaling) and to separate this general statement from the depth-$2$ notation used for the quantitative bound.

Q16 ("High level proof idea" opacity):

Thank you for this suggestion. In the camera-ready version, we will either remove this high-level paragraph from the main text or rewrite it so that it mirrors the formal proof structure in the appendix (dense-span argument, stationary condition, conditional-expectation step, and loss-optimality conclusion).

Q17 ("upstream $\times$ local" wording):

We will replace this by the standard phrasing "backpropagated (upstream) signal $\times$ local derivative", and we will include the explicit definition of the backprop signal. In our notation, the channel backprop signal is
\[
B_k^{(\ell)}(t;X,W)=\mathbb E_{C_\ell}\!\left[W_{C_\ell,k}^{(\ell)}\,\phi_\ell'(H_\ell(t,C_\ell;X,W))\,w_\ell(t,C_\ell)\right].
\]

Q18 (Meaning of $\partial_2 \mathcal{L}$):

Yes. We will define $\partial_2\mathcal L(y,\hat y)$ as the partial derivative with respect to the second argument (prediction).

Q19 (Meaning of "shifts" and heterogeneous shifts):

By shift we mean the additive translation behavior of intermediate-layer weights under i.i.d. initialization in the full-rank setting. In Nguyen \& Pham (2023) (Section 5.1.2, degeneracy discussion), intermediate-layer dynamics can collapse so that weights evolve essentially by a neuron-independent translation:
\[
w_i^\infty(t,C_{i-1},C_i)-w_i^\infty(0,C_{i-1},C_i)=w_i^\#(t),
\]
under constant initial biases and standard architectures. This is a precise identical-shift phenomenon.

In our RF-LR architecture, the frozen mixing matrix $L$ breaks this symmetry at the level of pre-activations:
\[
H_2(c_2;x,W)=\sum_{k=1}^r W_{c_2,k}f_k(x;W),
\]
so different neurons $c_2$ compute different linear combinations of channel partial functions. In the full-rank infinite-width limit (Nguyen \& Pham (2023), Theorem 54 and Corollary 30), intermediate layers can exhibit collapse: pre-activations become neuron-independent under i.i.d. initialization, and weights receive identical additive shifts (a neuron-independent translation). Here $W$ enforces heterogeneity across neurons by construction: each neuron $c_2$ computes a distinct channel combination $\sum_k W_{c_2,k}f_k$. Consequently, the backprop signal aggregated through $W_{C_2,k}$ (the $B_k$ term above) differs across neurons/channels, preventing the all-neurons-move-by-the-same-increment mode. This plays the same conceptual role as Nguyen \& Pham (2023)'s bidirectional diversity achieved through correlated initializations (Theorem 38), but realized architecturally rather than initializationally.

Q20 (Init index notation in Theorem~3.3):

Thank you for highlighting this ambiguity. Here, $\{n_1,n_2\}$ are exactly the numbers of neurons in the first and second hidden layers in the depth-2 quantitative theorem, and $\mathsf{Init}$ denotes a width-indexed family of initialization laws. This matches the neuronal-embedding setup used in Nguyen \& Pham (2023) (Section 4.1, Definition 1), where widths index the initialization family. We will rewrite the theorem in standard finite-width language (take widths $n_1,n_2$ and initialize i.i.d.) and remove the potentially confusing index-set wording in the main text. We will also state explicitly that the detailed quantitative proof is written for depth 2, while the underlying argument extends naturally to deeper architectures.

Q21 ("coupling procedure" and distance $D$):

We apologize for the confusing phrase "coupling procedure" in the submission: it was an editorial artifact and not intended to refer to a separate algorithm beyond the standard Nguyen \& Pham (2023) construction. What we mean is the usual finite-width initialization + embedding viewpoint: neurons are labeled by i.i.d. draws $C_i(j_i)$ from the limiting laws, and one compares empirical neuron averages to their mean-field expectations. There is no additional ad hoc coupling step beyond what is already implicit in the mean-field scaling.

For the distance denoted schematically by $D$ in the informal theorem statement, we will align notation with Nguyen \& Pham (2023) and with our manuscript appendix: the quantitative object is a supremum-norm-type discrepancy $D_T$ between finite and mean-field trajectories on $[0,T]$, and the final bound scales as $\mathcal O(1/\sqrt{n_{\min}}+\sqrt{\varepsilon})$ up to the usual exponential Gr\"onwall prefactor. When we refer to Wasserstein-type metrics in prose, we mean the transport viewpoint underlying those coupled $L^1$ integral assumptions, not a separate ad hoc metric introduced only for exposition.



Reviewer A6kZ (guarantees, high-$d$ mechanism, frozen-draw sensitivity)

(1) Convergence guarantees.\par
\noindent Question.

Convergence guarantees. Are there conditions under which convergence is guaranteed for RF-LR models (those mentioned in section~3.5 or otherwise)?

Answer. Our main guarantee is conditional on convergence of the mean-field dynamics. Unconditional convergence rates for deep RF-LR remain largely open; we will clarify this and avoid implying a general unconditional guarantee. We will also point to the sharpest partial results available and clarify which assumptions are sufficient for the global-optimality conclusion given convergence.

(2) High-dimensional extensions of the spike mechanism.\par
\noindent Question.

High-dimension extensions. Does the spike learning mechanism extend to higher dimensions with overlapping spikes in feature space?

Answer. Yes in substance: the mechanism (channel dominance $\leftrightarrow$ amplified backward signal through the mixing matrix and ReLU gates) is dimension-agnostic at the ODE algebra level because it depends on the same identity
\[
  H_2(c_2;x,W)=\sum_{j=1}^r W_{c_2,j}\,f_j(x;W)
\]
and the channelwise backprop objects that feed the mean-field equations. What we do not currently have for general input dimension $d$ is a clean theorem isolating well-separated (non-overlapping) spikes in feature space; overlapping channel contributions when $d>1$ complicate bookkeeping and the rigorous two-point caricature used in Theorem~4.1. We will state this distinction clearly in the camera-ready and add empirical evidence in higher $d$ where space permits.

Additional context --- Remark after Theorem~4.1 (mechanism beyond the two-point toy; drafted mainly for Reviewer~2 on Theorem~4.1, also cited here because the ODE algebra is dimension-agnostic): The explanatory content is not the simplified two-point, two-channel statement itself, but the positive feedback it isolates. Fix a second-layer neuron index $c_2$ and recall $H_2(c_2;x,W)=\sum_{j=1}^r W_{c_2,j} f_j(x;W)$. For ReLU, gating enters mean-field gradients through $\varphi'_2(H_2)=\mathbf{1}\{H_2>0\}$. The channel-$k$ backward signal (cf.~the ODE for $\partial_t w_1(\cdot,k)$) is built from expectations over $C_2$ of terms that couple (i)~the mixing weight $W_{C_2,k}$, (ii)~the gate $\mathbf{1}\{H_2(C_2;x,W)>0\}$, and (iii)~upstream factors (output error and later-layer weights). When channel $k$ dominates, neurons $c_2$ that contribute to the channel-$k$ drift are precisely those with the gate ``on.'' Across $c_2$, the effective weights in the expectation are then positively aligned with $W_{c_2,k}$ and with $\mathbf{1}\{H_2(c_2)>0\}$, producing a large contribution to the drift of the $k$-th channel, which in turn pushes $f_k$ further in the same direction: a reinforcing loop. The ODE-level mechanism is dimension-agnostic; the two-point theorem is a minimal tractable instance.

Visualization note. The figures emphasize $d=1$ mainly for clarity of visualization; the same feedback loop operates when $x\in\mathbb{R}^d$. Overlapping spikes mainly affect how we illustrate and formally isolate regimes, not the core mean-field structure above.

(3) Sensitivity to the frozen feature draw.\par
\noindent Question.

Sensitivity of randomness of fixed feature maps. How sensitive are the convergence guarantees and empirical performance to the initial draw of the frozen feature maps for finite width networks?

Answer. The analytical conditions are not tied to one sampling recipe. Once mixing matrices are frozen, the arguments use an almost-sure channel-aggregation bound (schematically $\|W^{(\ell)}\|_{\infty,1}\le rK$) and a richness/diversity assumption for the frozen first-layer features. In principle each frozen $W^{(\ell)}$ can be drawn from any law with finite support (or any law satisfying the same $\ell_{\infty,1}$ control a.s.); there is no Gaussian-specific requirement. In experiments we often use unbounded Gaussian initializations for simplicity at finite width; that is an engineering choice while the theorems use stability/richness predicates. Nonnegative / NMF-style $W\ge 0$ matches the same boundedness template and is intuition for some collapse modes; the ReLU dead-gate caveat is separate (high-probability-in-$r$ only as heuristic; cf.~Reviewer~VCqo Q10 and the ReLU appendix in the rebuttal PDF). At finite width, performance can vary with the draw. We have not run a systematic post-review sweep over independent frozen-feature draws with optimization held fixed; we will state this in the camera-ready, separate proved assumptions from mean-field idealization, and add compact multi-draw diagnostics if space permits or list them as future work.
